\clearpage
\VerdierPhysicalPageAnchor{0043}%
\VerdierSemanticAnchor{ch1-1-3-8-composition-cont}
$G^\circ$-圏 $(\mathcal C,G^\circ,T,c)$ に対応する次数付き圏 $\mathcal C[G]$
における合成の定義 (1.3.2.5) と両立することは直ちに確かめられる。

\VerdierSemanticAnchor{ch1-1-3-9}
\noindent\textbf{1.3.9.}\quad
$(\mathcal C,G,T,c)$ を $G$-圏とする。$(\mathcal C,G,T,c)$ の反対圏とは、
次のように定義される $G^\circ$-圏（$G^\circ$ は $G$ の反対群）
$(\mathcal C^\circ,G^\circ,T^\circ,c^\circ)$ である。圏
$\mathcal C^\circ$ は $\mathcal C$ の反対圏である。$g\in G^\circ$ に対する
関手 $T^\circ(g)$ を

\VerdierSemanticAnchor{ch1-1-3-9-eq1}
\begin{equation}
  \begin{aligned}
    T^\circ(g)X &= T(g^{-1})X,
      & X&\in\operatorname{Ob}(\mathcal C^\circ)
             =\operatorname{Ob}(\mathcal C),\\
    T^\circ(g)u &= T(g^{-1})u,
      & u&\in\operatorname{Fl}(\mathcal C^\circ)
             =\operatorname{Fl}(\mathcal C),
  \end{aligned}
  \tag{1.3.9.1}
\end{equation}
で定め、$\mathcal C^\circ$ における同型 $c^\circ(g,h)$ を

\VerdierSemanticAnchor{ch1-1-3-9-eq2}
\begin{equation}
  \begin{aligned}
    c^\circ(g,h)&\colon
      T^\circ(h\perp g)\longrightarrow T^\circ(h)T^\circ(g),\\
    c^\circ(g,h)&=c^{-1}(g^{-1},h^{-1})
  \end{aligned}
  \tag{1.3.9.2}
\end{equation}
で定める（$G^\circ$ における合成を $\perp$ で表す）。

\VerdierSemanticAnchor{ch1-1-3-10}
\noindent\textbf{注意 1.3.10.}\quad
$(\mathcal C,G,T,c)$ に対応する $G^\circ$ 型の次数付き圏を
$\mathcal C[G^\circ]\longrightarrow G^\circ$ とする。この圏は標準的切断系を
もつ。反対圏をとって得られるコファイバー圏を
$\mathcal C[G^\circ]^\circ\longrightarrow G$ とする。この圏は標準的余切断系を
もち、$G$ は群なので $G$ 上にもファイバー化されている。したがって
$\mathcal C[G^\circ]^\circ\longrightarrow G$ は標準的切断系をもち、そこから
$(\mathcal C,G,T,c)$ の反対圏と標準的に同型な $G^\circ$-圏が得られることは
直ちに確かめられる。

\VerdierSemanticAnchor{ch1-1-3-11}
\noindent\textbf{1.3.11.}\quad
$(\mathcal C',G',T',c')$ を $G'$-圏、
$(\mathcal C'',G'',T'',c'')$ を $G''$-圏とする。圏
$\mathcal C'\times\mathcal C''$ は明らかな仕方で標準的に
$G'\times G''$-圏の構造をもつ。

\VerdierSemanticAnchor{ch1-1-3-11-product}
得られた $G'\times G''$-圏を、$G'$-圏 $(\mathcal C',G',T',c')$ と
$G''$-圏 $(\mathcal C'',G'',T'',c'')$ の\emph{積}という。

\VerdierSemanticAnchor{ch1-1-3-12}
\noindent\textbf{1.3.12.}\quad
$(\mathcal C,G^\circ,T,c)$ を $G^\circ$-圏とする。実際には
$G^\circ$ の $\mathcal C$ への作用に右記法を用いることが最も多い。
$\mathcal C$ の各対象 $X$ に対し

\VerdierSemanticAnchor{ch1-1-3-12-eq1}
\begin{equation}
  T(g)X=X[g]
  \tag{1.3.12.1}
\end{equation}
とおく。同様に $\mathcal C$ の各射 $u$ に対し

\VerdierSemanticAnchor{ch1-1-3-12-eq2}
\begin{equation}
  T(g)u=u[g]
  \tag{1.3.12.2}
\end{equation}
とおく。したがって同型 $c(g,h)$ は各対象 $X$ に対し同型

\VerdierSemanticAnchor{ch1-1-3-12-eq3}
\begin{equation}
  c(g,h)(X)\colon X[gh]\longrightarrow X[g][h]
  \tag{1.3.12.3}
\end{equation}
を定める。

\VerdierSemanticAnchor{ch1-1-3-12-notation}
この記法には特に、$G^\circ$ の $\mathcal C$ への作用を $G$ の
$\mathcal C$ への右作用に変えるという利点がある。
